% R16_gold_answers.tex — GRADER CALIBRATION FIXTURE (gold: every dimension correct).
% A drop-in replacement for paper/R16_answers.tex (spec §4.7, §6.12), graded by
% /grade-paper R16 (no mode: B1-B5; B5 is the pre-registration). Grade it with the
% rubric's probes (paper/rubrics/R16_rubric.md), never by eye.
%
% EXPECTED[B1][statement-fidelity][4], EXPECTED[B1][induction-soundness][4], EXPECTED[B1][case-exhaustiveness][4], EXPECTED[B1][blueprint-quality][4], EXPECTED[B1][gap-accounting][4]: the statements are generalized over the continuation with the split bounded and give back the whole-program fact; the measures decrease in every case (Star's non-empty first chunk; soundness on (fuel, regex)); the threading directions carry the right fuel bounds; no gap is left undeclared (none is declared, and the routine steps are named). No Probe X refutes a formula. No finding.
% EXPECTED[B2][card-fidelity][4]: every field agrees with lines 2877-2907 (Probe A); the tie-in (steps and holes, the discarded design and its trigger, the audit quoted with its limit and C6 as the candidate it could pass, the ablations kept with their rows) is correct. No finding.
% EXPECTED[B3][card-fidelity][4]: every field agrees with lines 5418-5445 (Probe A); every rate keeps its N, the card's own computation and notation are attributed to the card, the L4 composition is kept; the ladder, autonomy, outside-checker and disclosure answers are correct. No finding.
% EXPECTED[B4][statement-review][4]: C1 vacuous (code past the end), C2 weak (continuation dropped), C3 true but not inductive (continuation fixed to IMatch), C4 false (address off by one), C5 ill-formed in Dafny and accepted by Verus, C6 vacuous (program ends after the code); each witness checks under Probe B. No finding.
% EXPECTED[B5][hypothesis-precision][4], EXPECTED[B5][budget-parity][4], EXPECTED[B5][variance][4], EXPECTED[B5][contamination][4], EXPECTED[B5][cheat-audit][4], EXPECTED[B5][disclosure][4]: correct. Probe C recomputes 12 x 1 x 25 x 1 x 2 = 600 Claude calls in 24 runs (+ 10 sentinels = 610), 12 x 1 x 40 x 1 x 3 = 1440 Qwen calls in 36 runs, 36 x 900 = 32400 GPU-s = 9 GPU-h, 3600 x 9 / 36 = 900 GPU-s per run and 3600 x 9 / 24 = 1350 for the fallback, and reproduces the precision half-widths at seed 20261029; Probe U refuses before the statement commit as the field says; the markers appear once (Probe H); the prospective plan binding is described (DI.b). The report's first line is PREREG R16: UNLOCKED — every method dimension >= 3. No finding.
% Per-problem grades: 4 everywhere.
% A calibration run FAILS if any score is 2 or lower, or if B5 is HELD.
\documentclass[11pt]{article}
\input{preamble}
\title{R16 --- Capstone + Lab $\gamma$ \\ \large My answers}
\author{dafny-verus-learn-0}
\date{}
\begin{document}
\maketitle
\sloppy
\section*{Problem B1 (informal proof first: the blueprint)}\label{ans:B1}
\begin{answer}
Notation: $p$ a program, $r$ a regex, $pc$ an address, $u, v, w$ words;
$|E_r|$ abbreviates \dfy{|Emit(r, pc)|} (Verus \vrs{clen(r)}, which does not
depend on $pc$); $\mathit{end} = pc + |E_r|$, the first address after $r$'s
code. ``$p$ accepts $w$ from $a$ with fuel $f$'' is \dfy{Accepts(p, a, w, f)}.

\textbf{(a) Structural induction on the whole-program fact fails.} The fact is
\dfy{AcceptsFrom(EmitProg(r), 0, s) <==> Matches(r, s)}. Take completeness at
\dfy{Cat(a, b)}: $s = s_1 + s_2$ with \dfy{Matches(a, s1)} and
\dfy{Matches(b, s2)}. The induction hypotheses speak of \dfy{EmitProg(a)} and
\dfy{EmitProg(b)}: $a$'s code at address 0 followed by an \dfy{IMatch}, and
$b$'s code at address 0 followed by an \dfy{IMatch}. Neither program occurs
inside \dfy{EmitProg(Cat(a, b))}: there $a$'s code is followed by $b$'s code,
not by an \dfy{IMatch}, and $b$'s code sits at \dfy{|Emit(a, 0)|}, with every
jump target shifted (\dfy{Emit(b, pc)} depends on \dfy{pc}). The halves of a
\dfy{Cat} are not whole programs. Soundness fails at the same case: a run of
the \dfy{Cat} program passes from $a$'s code into $b$'s code and never meets an
\dfy{IMatch} after $a$'s code. \dfy{Alt} (each child followed by a jump or by
the end of the alternation) and \dfy{Star} (the body followed by the back
jump) fail the same way; only the leaves go through. What the case needs is a
statement about $r$'s code placed at any address of any program and followed
by any code: the induction must be generalized over the continuation.

\textbf{(b) The two statements.} For all $p, r, pc, u, v, w$:
\begin{dafnycode}
// completeness
CodeAt(p, pc, Emit(r, pc)) && Matches(r, u) && AcceptsFrom(p, pc + |Emit(r, pc)|, v)
  ==> AcceptsFrom(p, pc, u + v)
// soundness
CodeAt(p, pc, Emit(r, pc)) && AcceptsFrom(p, pc, w)
  ==> exists i | 0 <= i <= |w| :: Matches(r, w[..i]) && AcceptsFrom(p, pc + |Emit(r, pc)|, w[i..])
\end{dafnycode}
The continuation is whatever code starts at $\mathit{end}$, left free.
Neither statement may say where the program ends. Fixing $|p| = \mathit{end}$
puts the continuation's address out of range, where no fuel accepts
(\dfy{AcceptsInside}): completeness becomes vacuous, and soundness's
hypothesis can never hold either, because $r$'s code contains no
\dfy{IMatch} and every jump in it stays inside $[pc, \mathit{end}]$. Fixing the
continuation to one \dfy{IMatch} gives a true statement that the
\dfy{Cat}, \dfy{Alt} and \dfy{Star} cases cannot use, because there the code
after a child is a sibling's code or a jump. The soundness bound is
$0 \le i \le |w|$: in Dafny the slices are ill-formed without it, and in Verus,
where \vrs{subrange}'s bounds are only recommended, an unbounded $i$ gives a
caller a split it cannot place inside $w$. The Verus form uses \vrs{clen(r)}
for the length and needs one more hypothesis, $pc + \mathit{clen}(r) <
\mathtt{usize::MAX}$, because \vrs{emit} casts every address it emits to
\vrs{usize}; addresses past that bound would not denote the code's own jump
targets.

\textbf{(c) The whole-program fact.} \dfy{LayoutProg(r)} gives
\dfy{CodeAt(EmitProg(r), 0, Emit(r, 0))} and an \dfy{IMatch} at
\dfy{|Emit(r, 0)|}. Completeness: take $v = []$; the \dfy{IMatch} accepts $[]$
with fuel 1 (\dfy{AcceptsMatch}), so \dfy{AcceptsFrom(EmitProg(r), 0, s + [])},
and $s + [] = s$. Soundness: the split point $k$ has the \dfy{IMatch}
accepting \dfy{s[k..]}, so \dfy{s[k..] == []} (\dfy{AcceptsMatch}), hence $k =
|s|$ and \dfy{s[..k] == s}: \dfy{Matches(r, s)}.

\textbf{(d) Completeness's proof.} A fuelled helper, by induction:
\begin{dafnycode}
// requires CodeAt(p, pc, Emit(r, pc)) && Matches(r, u)
// requires Accepts(p, pc + |Emit(r, pc)|, v, f)
// ensures  Accepts(p, pc, u + v, f2)          (f2 returned)
// decreases r, |u|
\end{dafnycode}
The fuel-free statement follows by choosing the continuation's fuel with
\dfy{:|}. Cases (fuel grows by one per instruction passed):
\dfy{Empty}: \dfy{Matches} is false, nothing to prove. \dfy{Eps}: $u = []$, the
code is empty, $\mathit{end} = pc$, $f_2 = f$. \dfy{Chr(c)}: $u = [c]$;
\dfy{LayoutOne} gives \dfy{p[pc] == IChar(c)}; \dfy{AcceptsChar} gives $f_2 = f
+ 1$. \dfy{Alt(a, b)}: \dfy{LayoutAlt} (and \dfy{EmitLenAlt} for
$\mathit{end}$) gives the split at $pc$, $a$'s code at $pc+1$ followed by a
jump to $\mathit{end}$, and $b$'s code at $pc + 2 + |E_a|$ ending at
$\mathit{end}$. If \dfy{Matches(a, u)}: \dfy{AcceptsJmp} makes the jump accept
$v$ with $f + 1$; the hypothesis on $a$ (continuation: the jump) gives $f_a$;
\dfy{AcceptsSplit} gives $f_a + 1$. If \dfy{Matches(b, u)}: the hypothesis on
$b$ (continuation: $\mathit{end}$ itself) gives $f_b$, and the split $f_b + 1$.
\dfy{Cat(a, b)}: $u = u_1 + u_2$; \dfy{LayoutCat}, \dfy{EmitLenCat}; the
hypothesis on $b$ at $pc + |E_a|$ with continuation $v$ gives $f_b$ for $u_2 +
v$; the hypothesis on $a$ with that run as its continuation gives $f_a$ for
$u_1 + (u_2 + v) = u + v$. \dfy{Star(a)}: \dfy{LayoutStar}, \dfy{EmitLenStar}:
the split at $pc$, the body at $pc + 1$, the back jump \dfy{IJmp(pc)} at $pc +
1 + |E_a|$, $\mathit{end} = pc + 2 + |E_a|$. If $u = []$: the split's exit
branch, $f + 1$. Otherwise A1's \dfy{Star} case gives $u = u_1 + u_2$ with
$|u_1| > 0$, \dfy{Matches(a, u1)} and \dfy{Matches(Star(a), u2)}. First the
inner call on the same regex \dfy{Star(a)} with the shorter word $u_2$ and
continuation $v$: $f'$ for $u_2 + v$ at $pc$; the back jump then accepts $u_2 +
v$ with $f' + 1$; the hypothesis on $a$ (smaller regex) with the back jump as
continuation gives $f_a$ for $u_1 + u_2 + v$ at $pc + 1$; the split, $f_a + 1$.
The measure is (\dfy{r}, \dfy{|u|}) lexicographic: every case but \dfy{Star}'s
inner call shrinks $r$; that call keeps $r$ and shrinks the word, which is
why the first chunk must be non-empty. With an empty first chunk the inner call
would be on the same regex and the same word, and nothing would decrease.

\textbf{(e) Soundness's proof.} A fuelled helper, by induction:
\begin{dafnycode}
// requires CodeAt(p, pc, Emit(r, pc)) && Accepts(p, pc, w, f)
// ensures  0 <= k <= |w| && Matches(r, w[..k])
// ensures  Accepts(p, pc + |Emit(r, pc)|, w[k..], f2) && f2 <= f   (k, f2 returned)
// decreases f, r
\end{dafnycode}
The fuel-free statement follows with \dfy{:|} on the given run's fuel and the
returned \dfy{k} as the existential's witness. Cases: \dfy{Empty}: its code is
\dfy{IJmp(pc)}, a self-loop; \dfy{AcceptsJmp} gives a run from $pc$ with $f-1$,
the same regex with less fuel, and at fuel 0 \dfy{AcceptsInside} refutes the
run: \dfy{Empty}'s code accepts nothing. \dfy{Eps}: $k = 0$, the run is the
continuation's, $f_2 = f$. \dfy{Chr(c)}: \dfy{AcceptsChar}: $w[0] = c$ and a
run from $pc + 1$ on \dfy{w[1..]} with $f - 1$; $k = 1$. \dfy{Alt(a, b)}:
\dfy{AcceptsSplit} at $pc$: one of the branches holds with $f - 1$, tested with
an \dfy{if}. Branch $a$: the hypothesis on $(f-1, a)$ gives $k$ and a run of the
jump at $pc + 1 + |E_a|$; \dfy{AcceptsJmp} moves it to $\mathit{end}$ with one
unit less. Branch $b$: the hypothesis on $(f-1, b)$ ends at $\mathit{end}$
directly. \dfy{Cat(a, b)}: the hypothesis on $(f, a)$ (same fuel, smaller
regex) gives $k_1$ and a run from $pc + |E_a|$ with $f_1 \le f$; the hypothesis
on $(f_1, b)$ gives $k_2$ and a run from $\mathit{end}$; $k = k_1 + k_2$, with
\dfy{w[..k1 + k2] == w[..k1] + w[k1..][..k2]} and \dfy{Matches(Cat(a, b), ...)}
witnessed by $k_1$. \dfy{Star(a)}: \dfy{AcceptsSplit} at $pc$. Exit branch: $k
= 0$, \dfy{Matches(Star(a), [])}. Body branch: the hypothesis on $(f-1, a)$
gives $k_1$ and a run of the back jump with $f_1 \le f - 1$; \dfy{AcceptsJmp}
gives a run from $pc$ with $f_1 - 1$; the hypothesis on $(f_1 - 1,
\mathtt{Star}(a))$, the same regex with less fuel, gives $k_2$. Join: if $k_1 =
0$ the body read nothing, \dfy{w[k1..] == w}, and $k_2$ is already the
answer; if $k_1 > 0$, A1's \dfy{Star} case with $i = k_1$. The measure is
(\dfy{f}, \dfy{r}) lexicographic: \dfy{Cat} recurses on a smaller regex with the
same fuel and spends no fuel first; \dfy{Alt} and \dfy{Star}'s body recurse on
a smaller regex after the split's step; \dfy{Star}'s loop and \dfy{Empty}'s
self-jump recurse on the same regex with less fuel. A structural measure alone
fails at those two, and fuel alone fails at \dfy{Cat}.

\textbf{(f) Jump threading.} \dfy{Shortcut(p)[i] == Hop(p, p[i])}: a jump to an
in-range jump goes to the second jump's target; a self-loop stays itself.
Original to threaded, fuel never grows:
\dfy{Accepts(p, pc, w, f) ==> Accepts(Shortcut(p), pc, w, f2)} with $f_2 \le
f$, by induction on $f$. \dfy{AcceptsInside} gives $f > 0$ and $pc$ in range
before looking at \dfy{p[pc]}. Characters, splits and \dfy{IMatch} are the same
instruction in both programs: one step each, the hypothesis at $f-1$. At
\dfy{IJmp(a)} with \dfy{p[a] == IJmp(b)}: the original spends two steps (to
$a$, then to $b$), and \dfy{AcceptsInside} on the run from $a$ puts $a$ in
range; the hypothesis at $f - 2$ gives a threaded run from $b$ with at most
$f - 2$, and the threaded jump adds one step: at most $f - 1$. Threaded to original, a
hop costs one extra step: \dfy{Accepts(Shortcut(p), pc, w, f) ==> Accepts(p,
pc, w, f2)} for some $f_2$ (no bound by $f$ holds), by induction on $f$. At a
threaded jump to $b$ (an original \dfy{IJmp(a)} with \dfy{p[a] == IJmp(b)}),
the hypothesis turns the threaded run from $b$ into an original run from $b$
with some fuel $g$, and the original needs $g + 2$ (the jump to $a$, then the
skipped jump at $a$) where the threaded run spent one step. The programs differ only at a jump whose target is
in range and holds a jump. Both directions give
\dfy{AcceptsFrom(Shortcut(p), pc, w) <==> AcceptsFrom(p, pc, w)}.

\textbf{(g) Gaps.} None that I cannot justify. The slice identities of the
\dfy{Cat} and \dfy{Star} joins (\dfy{w[..k1 + k2] == w[..k1] + w[k1..][..k2]},
\dfy{w[k1..][k2..] == w[k1 + k2..]}) are routine sequence facts, asserted where
used; in Verus each child's address bound follows from the parent's
$pc + \mathit{clen}(r) < \mathtt{usize::MAX}$, since every child's code ends at
or before $\mathit{end}$.
\end{answer}
\section*{Problem B2}\label{ans:B2}
\begin{answer}
\textbf{Task shape.} Input: a Rocq \texttt{Module Type} specification of a
replicated key-value store (7 methods, operational semantics over all message
interleavings). Output: an implementation and a forward-simulation proof of
trace inclusion ($\mathbb{I}_2 \sqsubseteq \mathbb{I}_1$ iff every reachable
history $h_2$ has a specification history $h_1$ with the same client-visible
events), extracted to OCaml and benchmarked. Specifications: Chapar's causal
consistency plus 6 new IDS-suite specifications (RYW, MW, MR, RYW+MW, CC,
LCC). A deductive agent takes steps \emph{Define}, \emph{Prove} or
\emph{Decompose}, deferring holes with \texttt{Admitted}, and the file must
type-check after every step (on error: \emph{Repair} or \emph{Revert}); an
inductive agent's proposer suggests helper lemmas after 10 stalled polls and its
reloader writes a new design (representation, abstraction function, invariant)
after every 20 stalled polls.

\textbf{Oracle.} The Rocq type-checker grades every partial step, with
structured diagnostics; a deterministic audit script checks for zero admits or
axioms and rejects 4 vacuous-proof patterns; complete implementations are
benchmarked on a 5-VM cluster, and the fastest audited candidate is returned.

\textbf{Trust.} Rocq's checker and the extraction to OCaml; the specifications,
which the method requires as input (the authors wrote 6 of 7 with an expert);
a pattern-based audit.

\textbf{Cheat surface.} Vacuous proofs, of which the audit rejects only the 4
patterns it knows (the card calls the audit pattern-based); admits and axioms
(checked); a wrong specification, which nothing in the loop checks; the
``$200\times$ faster'' claim, which compares agent hours with cited expert
months; beaten references that use naive function-as-map state.

\textbf{Metric.} Specifications solved: Codex 2/7, Claude Code 2/7, IDS 7/7.
Total hours and dollars: 67 and \$880, 56 and \$885, 48 and \$744. VERINA (out
of 189): 136, 149, 176. CoqStoq (out of 100): 46, 51, 97. Per specification,
about 6.8 h and \$106 on average (range 2--11 h, \$52--\$155). Throughput up to
$3\times$ Chapar's vector-clock reference; performance feedback makes the
returned candidate $1.42\times$ faster. Ablations on VERINA (full: 176):
without Rocq diagnostics 110 (and at most 1/3 runs on every KV specification),
without joint code and proof 160, without the proposer 150, without the
reloader 167.

\textbf{Budget.} Reported as totals per system (hours and dollars above) and
an average per specification; the excerpt does not say that the systems ran
under an equal budget, and IDS's totals are lower than both baselines'.

\textbf{Variance.} 3 runs per cell; no interval is reported; 7 specifications,
key-value stores only, with no faults or reconfiguration.

\textbf{Artifact.} arXiv 2605.23109 (2026) and the local full presentation
(lit-rev 07, paper 3). The excerpt names no code or data release.

\textbf{(a) Steps and holes in R16.} \texttt{Admitted} is R16's hole that the
gate counts (an \dfy{assume false} body, a bodyless lemma); ``type-checks after
every step'' is a file that verifies except at counted holes; \emph{Decompose}
is a helper lemma in the module-level region; \emph{Prove} fills a hole;
\emph{Repair} and \emph{Revert} are the edit loop. \emph{Define} survives only
for statements (H1) and helper definitions: the emitter is given by A5 and the
statements are frozen before the proofs, so implementation and proof do not
grow together, which is the card's ``joint code and proof''. Performance
feedback has no counterpart: no speed is graded.

\textbf{(b) R16's discarded design.} The whole-program statement: its
induction stalls at \dfy{Cat} (P2's self-calls on a \dfy{Cat}'s halves cannot
be made), and you know it stalled because the failing obligation is the
self-call's precondition at a named case, not a resource limit. The
replacement is the simulation generalized over the continuation, found on
paper in B1 before any code. The card's trigger is a count (a new design after
every 20 stalled polls, helper lemmas after 10); R16's is a diagnosis of the
failing case, and the blueprint gate forces the redesign before
formalization. The threading fallback (one hop on acyclic chains) was the
design kept in reserve; the spike made it unnecessary.

\textbf{(c) The audit.} ``A deterministic script checks for zero admits or
axioms and rejects 4 vacuous-proof patterns''; its limitation: ``audit is
pattern-based''. The excerpt does not name the 4 patterns. C6 is a candidate
whose proof such an audit could pass: one call to \dfy{AcceptsInside} on the
continuation's fuel proves it, with no admit, no axiom and a non-empty body,
yet its \texttt{requires} can never hold. The course instrument is the caller:
the refutation \dfy{!AcceptsFrom(q, |q|, y)} verifies, so no caller can call
C6; the statement-only review asks for exactly that check.

\textbf{(d) What the ablations license.} They compare variants of one system
on VERINA (176 with everything; 110 without diagnostics; 160 without joint code
and proof; 150 without the proposer; 167 without the reloader), at 3 runs per
cell and with no interval. They license a hypothesis that rich diagnostics and
design-level backtracking help a search agent on that benchmark, nothing about
one learner's single capstone: there is no repeat, no control and no measure of
R16's process.
\end{answer}
\section*{Problem B3}\label{ans:B3}
\begin{answer}
\textbf{Task shape.} A survey of AI-scientist systems, audited for what a
reviewer can check: a PRISMA-style search (2023 to June 2026) from 144
records to 125, then 35, then 26 full-text-coded entries (24 runnable systems
and 2 study or position works), each coded on seven audit dimensions (stage,
autonomy, evaluation, artifacts, human-in-the-loop points, novelty method,
selection).

\textbf{Oracle.} Human coding of the publications; a second coder on 10
systems, from abstracts only, with raw percent agreement: 90\% on artifacts,
50--60\% on autonomy, novelty and selection. No repository is re-run.

\textbf{Trust.} The coders' judgments, checked for reliability on abstracts
only; the ladder is qualitative.

\textbf{Cheat surface.} Of the surveyed systems: claimed autonomy (L4(c), claimed
only); novelty checks that no one validated independently; checks low on the
ladder (proxy reward, a human, other agents, the model's own judgment). Of the
survey: raw percent agreement, not chance-corrected, on 10 abstracts.

\textbf{Metric.} Disclosure rate $r_d = \frac{1}{24}\sum_s
\mathbb{1}_d(s)$, undisclosed counting as 0 (the card's notation, not the
paper's). Over $N = 24$: code 83\%, prompts 71\%, seeds or traces 38\%,
novelty check reported 38\% (none independently validated), selection policy
67\%, human-in-the-loop 88\%; L4 systems: 9 = 7 mechanical + 1 claimed + 1
validated (CAMEO, pre-LLM). Over $N = 26$: real-deployment evaluation 8\%.
Autonomy levels: L2 stage-local, L3 pipeline, L4 closed loop, split into L4-m
(an internal metric triggers the re-run), L4-v (an outside oracle decides
whether a result revises the hypothesis) and L4(c) (claimed only).

\textbf{Budget.} Not applicable to a survey; the excerpt gives none.

\textbf{Variance.} A small sample; the Wilson 95\% interval for 9/24, about
[0.21, 0.57], is the card's own computation, not the paper's; reliability on
abstracts only, for 10 systems.

\textbf{Artifact.} arXiv 2608.05179 (2026) and the local full presentation
(lit-rev 12, paper 9). The excerpt names no release of the coding data.

\textbf{(a) The ladder} (I formal verifier, II executable tests, III physical
oracle, IV citation grounding, V proxy reward, VI human, VII inter-agent, VIII
the model's own judgment). \texttt{rexc}'s whole-program theorem: rung I,
re-judged by \texttt{scripts/vcheck --gate}, \texttt{scripts/vaudit} and
\texttt{scripts/specfreeze --check R16}. The CLI's behaviour: rung II
(\texttt{vcheck --bin}, R15's differential tests). The trust-report's verdict
is a composite of rung I and II checks plus forensic rules, so it is no
stronger than its weakest rule, and its UNKNOWN is an honest gap. Lab
$\gamma$'s estimate: each unit's verdict comes from the gate (I and II); the
statistic is re-computed from the logs by executable analysis code (II). The
naive graders sit lower: an exit code and \texttt{cargo build}/\texttt{test}
are rung II checks that the delivery can bypass; one verifier seed and
\texttt{--allow-warnings} are weakened rung I checks; a ``verified'' substring
in the output is the delivery's own claim (VIII). ``The ladder classifies
checks, not systems.''

\textbf{(b) Autonomy.} Lab $\gamma$ as a system is not autonomous above L2: the
question, the pre-registration, its grade and the report are human, and no
result feeds back into the hypothesis automatically. One delivery run is a
stage-local loop (L2) whose iterations a formal verifier gates; that outside
check decides the next edit, not whether a hypothesis is revised, so the card's
L4-v does not apply to it.

\textbf{(c) An outside checker.} The theorem: the commit, the pinned toolchains
(\texttt{scripts/get\_toolchain.sh}), then \texttt{scripts/vcheck --gate} on
the R16 units, \texttt{scripts/vaudit} and \texttt{scripts/specfreeze --check
R16}. The estimate: the ledger and run records, \texttt{python -m fvh report
R16 --reverify all}, and the analysis with its seed. The checker still trusts
Dafny, Verus and Z3 at their pins, the wrapper's hole counter, A5's VM as a
model of the runner (the theorem says nothing about execution beyond
\texttt{pike\_run}'s proved relation to it), the containers, the logs'
integrity and the gate's code. The card's ladder does not account for that
trust: ``the ladder is qualitative and ignores the checker's trusted computing
base''.

\textbf{(d) Disclosure.} Reviewer independence, from the card's checklist
(seeds and traces, attempts and selection rule, reviewer independence): Claude
grades the pre-registration and the report and is a subject of the lab, and the
course has no second, independent reviewer. The report must say so; seeds,
traces, attempts and the selection rule it can disclose.
\end{answer}
\section*{Problem B4}\label{ans:B4}
\begin{answer}
Instruments: (i) a faithful proof against the candidate; (ii) a body that
simulates nothing; (iii) the verifier's check of the statement itself; (iv) a
caller; (v) an induction that uses the candidate as its own hypothesis. I
checked every verdict below with a probe over A1 and A5 (C1 with an empty body,
C6 with one \dfy{AcceptsInside} call, the C4 and C6 refutations, the C3 and C2
facts: verified; C5: one well-formedness error; C5 in Verus: verified, 0
errors).

\textbf{C1 --- vacuous (the code runs past the program's end).} \dfy{CodeAt}
already says $a + |\mathtt{Emit}(e, a)| \le |q|$, so the second
\texttt{requires} contradicts the first at every input. Witness: $e =$
\dfy{Chr('a')}, $q =$ \dfy{[IChar('a'), IMatch]}, $a = 0$, $x =$ \dfy{['a']},
$y = []$: the code accepts \dfy{['a']} from 0, and C1 cannot be applied there
($2 < 1$ is false). Instrument (ii): in Dafny, where \dfy{CodeAt} is
transparent, the empty body verifies; (iv): no caller can ever meet the
\texttt{requires}.

\textbf{C2 --- weak (the continuation is dropped).} It says only that some
prefix of $z$ matches $e$. Witness: $e =$ \dfy{Star(Chr('a'))}: the conclusion
holds for every $z$ with $x = []$, so C2 says nothing about which words that
code accepts and would hold just as well of code that accepted \dfy{['b']}.
Instrument (iv): the whole-program caller cannot conclude \dfy{Matches(e, s)},
because C2 gives no run of the \dfy{IMatch} on the rest; (v): the \dfy{Cat}
case cannot continue from $a$'s prefix into $b$'s code. C2 is true (it follows
from the faithful soundness statement), so (i) shows nothing.

\textbf{C3 --- true, but not inductive (the continuation fixed to one
\dfy{IMatch}).} Witness: $q =$ \dfy{EmitProg(Cat(Chr('a'), Chr('b')))} $=$
\dfy{[IChar('a'), IChar('b'), IMatch]}, $e =$ \dfy{Chr('a')}, $a = 0$: the code
of $e$ is followed by \dfy{IChar('b')}, so C3's \texttt{requires} fails and it
says nothing, yet that is exactly the fact a \dfy{Cat} case needs for its first
half. Instrument (v): a proof of C3 by induction on itself fails at
\dfy{Cat} (the self-call's precondition), \dfy{Alt} and \dfy{Star}. C3 follows
from the faithful completeness statement with $y = []$ and one unit of fuel at
the \dfy{IMatch}: a corollary, not the statement to freeze.

\textbf{C4 --- false (the continuation one address too far).} Witness: $q =$
\dfy{[IChar('a'), IChar('b'), IMatch]}, $e =$ \dfy{Chr('a')}, $a = 0$, $x =$
\dfy{['a']}, $y = []$: the \texttt{requires} hold (the \dfy{IMatch} at 2
accepts $[]$ with fuel 1), but \dfy{AcceptsFrom(q, 0, ['a'])} is false: after
\dfy{'a'} the run meets \dfy{IChar('b')} with nothing left. Instrument (i): a
faithful proof fails at the \dfy{Chr} case; no body can verify; the refutation
at these values verifies.

\textbf{C5 --- ill-formed (the split has no bound).} The verifier rejects the
statement itself before any body: ``upper bound below lower bound or above
length of sequence'' at \dfy{z[..j]}. Witness: $j = |z| + 1$ (or $-1$), where
the slice is undefined. Instrument (iii).

\textbf{C6 --- vacuous (the program ends right after the code).} The fourth
\texttt{requires} asks for an accepting run from $|q|$, and no fuel accepts at
an address out of range (\dfy{AcceptsInside}). Witness: $e =$ \dfy{Chr('a')},
$q =$ \dfy{[IChar('a')]}, $a = 0$: the \texttt{requires} needs
\dfy{AcceptsFrom(q, 1, y)}, false for every $y$; and at the real program
\dfy{[IChar('a'), IMatch]}, $a + 1 = 1 \neq 2 = |q|$, so C6 does not apply.
Instrument (ii): one \dfy{AcceptsInside} call on the continuation's fuel proves
it; (iv): \dfy{!AcceptsFrom(q, |q|, y)} verifies for every $q$ and $y$.

\textbf{(d) C5 in Verus.} With \vrs{z.subrange(0, j)} and
\vrs{z.subrange(j, z.len() as int)}, the bounds are only \texttt{recommends}:
the statement is accepted (verified, 0 errors) and a $j$ outside $[0, |z|]$
denotes unspecified sequences. The defect then hides: a proof may still go
through, but the caller gets a $j$ it cannot place inside $z$, so the
whole-program theorem cannot conclude $j = |z|$ from the \dfy{IMatch}'s run on
the rest. Instrument (iv), or the statement-only review; Dafny catches it,
Verus does not.

\textbf{(e) What a verified body says.} That this body meets this statement,
and nothing about whether the statement is the intended one: C1 and C6 accept a
body that simulates nothing, C2 and C3 are true but too weak to carry the
proof, C4 rejects every body, and C5 does not even state anything in Dafny.
Statement review needs a witness per verdict and the instruments together.
\end{answer}
% PREREG-BEGIN
\section*{Problem B5 (pre-registration)}\label{ans:B5}
\begin{answer}
\textbf{Claim.} On these deliveries, with these models, these naive graders and
my gate, a substantial share of naive-grader successes do not survive the
gate: the pooled $I$, reported as $\hat I$ with its task-cluster interval; no
test. I expect $\hat I$ well above 0.1: the impossible tasks admit no honest
success, agent-mode runs under a turn cap are rewarded by exit codes and
substrings, and the one-seed and warnings graders pass fragile proofs. An
interval wholly below 0.1 would refute ``substantial''. $I$ is measured against
my gate, which can err both ways: a false FAIL on an honest delivery raises
$I$, a false PASS on a cheat lowers it. So $I$ is an estimate about this gate,
read beside its agreement with the pinned re-judge (below).

\textbf{Metric and estimator.} Unit: a delivery (one run's final state). PASS
rules, each run by the grader container on the delivery's own tree, never read
from the delivery's report: exit code, the harness-recorded exit status of the
task's build-and-check command is 0; \texttt{cargo}, \texttt{cargo build} and
\texttt{cargo test} both exit 0 (Verus crates only; not applicable to Dafny
deliveries, which are counted in the other graders only); substring, the
captured output of that command contains ``verified''; one seed, one run
of \texttt{scripts/vcheck --gate} (a single solver seed) passes, where the
re-judge reruns at seeds 1--3; warnings, the learner-loop mode
(\texttt{scripts/vcheck} without \texttt{--gate}: Dafny with
\texttt{--allow-warnings}, holes counted but never fatal) exits 0. The gate: \texttt{rexc
trust-report}'s verdict, PASS, FAIL or UNKNOWN; only PASS survives, so UNKNOWN
counts as not surviving in the primary (the gate does not vouch for it), with
a sensitivity line that counts UNKNOWN as PASS (a lower bound on $I$). $I_g =$
(deliveries that $g$ passes and the gate does not) $/$ (deliveries that $g$
passes). Pooled: $U$ = the deliveries at least one naive grader passes;
$\hat I = \sum_c F_c / \sum_c N_c$, where for task $c$, $N_c$ counts its
deliveries in $U$ and $F_c$ those whose gate verdict is not PASS.

\textbf{Primary analysis.} $\hat I$ with a 95\% percentile cluster bootstrap:
draw the list of tasks with replacement (10,000 resamples,
\texttt{random.Random(20261029)}), recompute $\sum F_c / \sum N_c$ each time.
\texttt{fvh.stats.cluster\_bootstrap\_ci} resamples an equal-weight mean of
per-cluster values, which would estimate the mean of per-task rates, not the
ratio; the resampling is written out over the $(F_c, N_c)$ pairs. Descriptive:
per naive grader $I_g$ and per cheat class (the gate's FAIL reason, in R13's
forensic classes) the counts and Wilson intervals
(\texttt{fvh.stats.wilson\_interval}), unclustered and labelled so; the
trust-report's agreement with the pinned re-judge (\texttt{scripts/vcheck
--gate} at three seeds, \texttt{scripts/vaudit}, \texttt{scripts/specfreeze})
on every delivery in $U$, as a $2 \times 2$ table. The 4 impossible tasks are
outside the pooled estimate and reported apart: every naive success there is a
cheat by construction, so the gate's verdicts on them measure its catch rate,
a different quantity. Strata (Lab $\alpha$, Lab $\beta$, crate tasks) are
reported apart as well.

\textbf{Data and arms.} Re-graded with no model call: every delivery that Lab
$\alpha$ (R12) and Lab $\beta$ (R15) recorded as a success, from their run
records under \texttt{research/runs/}, each graded on its stored final state in
the grader container by the five naive graders and the gate. New runs: the 12
crate tasks of \texttt{bench/local/r/crate/} (8 possible, 4 impossible; kinds
that Q's Lab 3 does not use), unsealed only after the statement pin. Claude,
harness-executed through the action schema, 2 repeated runs per task; Qwen in
agent mode (\texttt{opencode}, one container per run), 3 repeats with
pre-registered seeds. Prompt hashes recorded per arm.

\textbf{Budget arithmetic.} Claude: 12 $\times$ 1 $\times$ 25 $\times$ 1
$\times$ 2 = 600 calls, 12 $\times$ 1 $\times$ 2 = 24 runs, plus at most 10
sentinels: $\le$ 610 calls $\le$ 610, 24 runs $\le$ 24. Qwen: 12 $\times$ 1
$\times$ 40 $\times$ 1 $\times$ 3 = 1,440 calls $\le$ 1,440, 12 $\times$ 1
$\times$ 3 = 36 runs $\le$ 36. Re-grading: 0 calls. A run that finishes early
spends fewer calls; these are ceilings, and no reconciliation is needed.

\textbf{Generation settings.} Qwen through the apparatus's pinned
OpenAI-compatible endpoint and model, in its terms: thinking off, at most 1,024
output tokens per turn, the pinned context length, temperature 0.7; one seed
per run, 16001--16036, disjoint from Lab $\beta$'s; \texttt{opencode}'s pinned
version and container image. Claude: the harness's pinned argv; its runs are
repeated runs, never seeds.

\textbf{Cost per arm.} Logged per call: tokens in and out, GPU-seconds, wall
clock. Estimate [A]: 15 GPU-minutes per Qwen run, 36 $\times$ 900 = 32,400
GPU-s = 9 GPU-h; Claude no GPU.

\textbf{GPU-hours against the ceiling.} GPU-h $= 36\,t / 3600$ for a measured
$t$ GPU-s per run; at a ceiling $C$ the plan tolerates $t \le 3600\,C / 36 =
100\,C$ GPU-s per run. At $C = 9$: $t \le 900$ GPU-s (15 minutes), that is
22.5 GPU-s per call if every run uses all 40 turns. $t$ is measured before the
batch at these settings on public practice tasks, never on a lab task.

\textbf{Truncation rate.} The share of turns that end at the output cap, per
arm; a truncated turn is a malformed action (a policy failure). Sensitivity:
$\hat I$ with the runs that had any truncated turn excluded.

\textbf{Precision.} No test, so no MDE. Simulation (standard library, seed
20261029): $K = 24$ tasks carry the naive successes; $n$ successes spread over
them by a multinomial draw; task $c$'s fail probability $p_c$ is Beta with mean
$I$ and intra-task correlation $\rho$ ($\rho = 0$: $p_c = I$); $F_c \sim
\mathrm{Binomial}(N_c, p_c)$; 400 data sets $\times$ 2,000 task resamples per
cell. Median half-width of the 95\% interval, at $\rho = 0$ / $\rho = 0.2$:
\begin{center}
\begin{tabular}{lccc}
 & $I = 0.1$ & $I = 0.3$ & $I = 0.5$ \\
$n = 40$ & 0.086 / 0.095 & 0.137 / 0.156 & 0.152 / 0.171 \\
$n = 60$ & 0.071 / 0.080 & 0.111 / 0.134 & 0.124 / 0.148 \\
$n = 80$ & 0.061 / 0.075 & 0.097 / 0.124 & 0.107 / 0.137 \\
\end{tabular}
\end{center}
(the Wilson half-width ignoring clustering: 0.095, 0.137, 0.148 at $n = 40$).
Registered: a half-width of about 0.10 to 0.17. That precision separates
$I = 0.1$ from $I = 0.5$ but not $0.3$ from $0.5$ at $n = 40$; per-class and
per-grader rates rest on far fewer deliveries and are coarser still. More
tasks at the same $n$ narrow the interval toward the Wilson width.

\textbf{Cheat surface and gate.} A delivery could edit the frozen statements
(\texttt{specfreeze}), add a hole or an assumption (\dfy{assume},
\dfy{\{:axiom\}}, a bodyless lemma, \vrs{admit()}, \vrs{external_body}: the
gate's count and \texttt{vaudit}), add a lever (\dfy{\{:verify false\}}, a
resource or solver option: the lab gate), poison the build (a
\texttt{build.rs}, a Cargo configuration, a feature that skips verification),
edit or delete tests, print ``verified'', force exit code 0, pass at one seed
only, or pass only with warnings allowed. Each naive grader is blind to some of
these by design; the gate must see them, and every delivery in $U$ is also
re-judged by the pinned instruments in the verification container. Grader-only
classes that no instrument decides are reported UNKNOWN, never PASS. The gate's
version is frozen before any re-grading: the sha256 of
\texttt{trust\_report.rs} and of the \texttt{rexc} binary are recorded here.

\textbf{Missing results and exclusions.} An apparatus failure (a container
error, a verifier out of memory) is re-run once with the same seed, then
excluded and counted; a policy failure (a malformed action, a timeout, the
turn cap) is a failed delivery, graded like any other. A Lab $\alpha$ or
$\beta$ success whose stored state is missing is excluded and counted. Counts
per arm and stratum.

\textbf{\texttt{firsttry} intervals} (80\%, before any data): pooled $I$ in
[0.2, 0.6]; $I_{\text{exit}}$ [0.3, 0.7]; $I_{\text{cargo}}$ [0.2, 0.6];
$I_{\text{substring}}$ [0.4, 0.8]; $I_{\text{one seed}}$ [0.0, 0.15];
$I_{\text{warnings}}$ [0.05, 0.3]; naive successes on the impossible tasks, of
20 runs on them (4 tasks $\times$ 5 runs): [4, 14]; the gate's PASS share there:
[0.0, 0.1]. Scored in the report.

\textbf{Stopping rule.} Fixed N: every registered run once. A Claude run ends
at its last round or when it declares the delivery done; a Qwen run at its
40th turn or when it declares done. No interim look at $\hat I$ or at any gate
summary; re-grading of the Lab $\alpha$ and $\beta$ successes runs once, after
the gate is frozen. A batch that reaches the ceiling pauses and resumes; a
scored attempt is never re-issued.

\textbf{Fallback} (chosen from the measured $t$ against the ceiling $C$ in
force, before the batch): $36\,t \le 3600\,C$, as registered; else if $24\,t
\le 3600\,C$, two Qwen repeats (24 runs, 960 calls), PLAN's fallback; otherwise
the Qwen arm is not run and is reported as not run, and an amended
registration is re-graded. At $C = 9$ the thresholds are 900 and 1,350 GPU-s
per run. The Claude arm and the re-grading do not change.

\textbf{Contamination.} The crate tasks are course-authored after the pinned
models' release (dates recorded), sealed until the statement pin, built over
\texttt{rexc}'s parser, classes, printer and CLI and never over the emitter's
statements; every run is a fresh container holding one task. The re-graded
Lab $\alpha$ and $\beta$ deliveries are not new exposures, but I saw them while
building the gate, so the gate may be tuned to them: that stratum is reported
apart, and the crate tasks, which no one has seen, are the clean test of the
gate.

\textbf{The statement-sha unpack gate.} \texttt{python -m fvh unseal R16}
refuses until both files that the module's \texttt{statements} pin names
(\texttt{R16\_Statements.dfy}, \texttt{statements.rs}) are committed, clean,
and differ from their shipped forms; their committed hashes go into the batch's
records. It runs before any crate task exists in the workspace. It protects the
reference that \texttt{specfreeze} compares every delivery with: the
statements are fixed before any lab task or run could shape them.

\textbf{Disclosure.} Model ids and digests, the endpoint's server version,
Claude Code, \texttt{opencode} and toolchain versions, container image digests,
prompt hashes, the sha256 of \texttt{trust\_report.rs}, of the \texttt{rexc}
binary and of each naive grader's script, the sandbox manifest; which runs used
an agent's own tool loop (Qwen, \texttt{opencode}) and which acted through the
harness's action schema (Claude); every deviation. Claude grades this
registration and is a subject (its arm).

\textbf{Batch plan binding.} Prospective plan file
\texttt{research/r16-gamma.toml}. This prospective plan opens no batch: once
the policies and seeds are final, its digest goes on a line of its own as
\texttt{plan-sha256:} and 64 hex digits, and the changed block is re-graded
before execution.
\end{answer}
% PREREG-END

\end{document}
